\par\noindent\textbf{(1)} 根据van Cittert–Zernike定理，归一复相干度 \[\gamma_{12}=\int I(\boldsymbol\theta)e^{-i2\pi\boldsymbol B\cdot\boldsymbol\theta/\lambda}d^2\theta/\int I(\boldsymbol\theta)d^2\theta\]。均匀圆光源 \(I(\theta)=I_0\) 在 \(\theta\le\theta_s\)，所以二维积分为圆盘Fourier变换。
\par\noindent\textbf{(2)} 取基线沿极坐标零方向，角积分是 \(2\pi J_0(2\pi B\theta/\lambda)\)。径向积分 \(\int_0^{\theta_s}\theta J_0(q\theta)d\theta=\theta_sJ_1(q\theta_s)/q\)，归一后 \(\gamma_{12}=2J_1(x)/x\)，其中 \(x=2\pi B\theta_s/\lambda\)。
\par\noindent\textbf{(3)} 第一非零 \(J_1\) 根为 \(x_1\approx3.831706\)。因此 \(B_1=x_1\lambda/(2\pi\theta_s)\approx0.610\lambda/\theta_s\)。若用圆盘角直径 \(\Theta=2\theta_s\) 表示，则 \(B_1\approx1.22\lambda/\Theta\)，可以由干涉可见度零点反推光源角尺寸。
\par\noindent\textbf{(4)} 若真实光源存在临边昏暗、热点或不对称结构，角强度不再是均匀圆盘，\(\gamma_{12}\) 需对实际亮度作归一二维Fourier变换，零点位置和相位均可能变化。该定理要求源点之间统计不相干，不适用于单一空间相干激光圆孔的相同解释。
